% !PRINTED 136
\weberpage{136}
%<<CONT

des grössten gemeinsamen Theilers auf diese beiden Functionen anwendet, so erhält man die Bedingung gemeinsamer Factoren daraus, dass man den letzten constanten Rest gleich Null setzt in Gestalt einer Gleichung zwischen den Coëfficienten:
\[
 F(a_1,a_2,\ldots,a_n)=0,
\]
worin $F$ eine ganze rationale Function der Argumente $a_1,a_2,
\ldots,a_n$ ist, die die \gesperrt{Discriminante} von $f$ heisst, deren Eigenschaften und Bildungsweise im nächsten Abschnitt noch eingehender behandelt werden soll; jedenfalls kann, da es überhaupt Gleichungen $n$ten Grades ohne mehrfache Wurzeln giebt, $F$ nicht identisch verschwinden.

Betrachten wir also jetzt
\[
 F(a_1+\varepsilon_1,a_2+\varepsilon_2,
\ldots,a_n+\varepsilon_n),
\]
so wird auch diese Function, wenn die $a_1,a_2,
\ldots,a_n$ bestimmte, die $\varepsilon_1,\varepsilon_2,
\ldots,\varepsilon_n$ aber unbestimmte Zahlen sind, nicht identisch verschwinden, und man kann über die $\varepsilon$ so verfügen, dass ihr Werth von Null verschieden ausfällt, auch dann noch, wenn für die absoluten Werthe der $\varepsilon$ eine beliebige obere Grenze festgesetzt ist.\footnote[1]{Dies wird erst später streng begründet werden. Wir setzen es bei diesen vorläufigen Bemerkungen, auf die weiter keine Schlüsse gegründet werden, einstweilen voraus.}

Man sieht also aus diesen Betrachtungen, wie eine $a$-fache Wurzel $\alpha$ von $f(x)$ bei stetiger Veränderung der Coëfficienten in $a$ einfache Wurzeln auseinander strahlt, so dass man auch von diesem Gesichtspunkte berechtigt ist, die $a$-fache Wurzel als durch das Zusammenfallen von $a$ einfachen Wurzeln entstanden zu betrachten.

Nehmen wir den Coëfficienten der höchsten Potenz von $x$ nicht gleich $1$ an, untersuchen also die Gleichung in der Form
\begin{equation}
 a_0x^n+a_1x^{n-1}+a_2x^{n-2}+\cdots+a_{n-1}x+a_n=0,
\tag{21}
\end{equation}
so sind, wenn $a_n=0$ ist, eine oder mehrere der Wurzeln gleich Null, und die Gleichung (21) lässt sich durch $x$ oder eine höhere Potenz von $x$ dividiren, wodurch eine Gleichung von entsprechend niedrigerem Grade entsteht, in der das letzte Glied von Null verschieden ist. Wir wollen also von vornherein $a_n$ von Null verschieden annehmen und nun in (21)
\begin{equation}
 x=\frac1y
\tag{22}
\end{equation}
%CONT>>
